\documentclass[11pt]{article}
\usepackage[margin=1in]{geometry}
\usepackage{amsmath,amssymb}
\usepackage{booktabs}
\usepackage{hyperref}
\hypersetup{colorlinks=true,linkcolor=blue,urlcolor=blue}

\newcommand{\R}{\mathbb{R}}
\newcommand{\HSIC}{\mathrm{HSIC}}

\title{\textbf{Discovering the Metric Linkage with HSIC\\
\large a minimal prototype}}
\author{}
\date{}

\begin{document}
\maketitle

\section{Idea}

Business metrics live at different levels (order, user, merchant) and have
\emph{no explicit formula} linking them. We only want to know \emph{which metrics
move together} --- the association \emph{linkage} --- not any causal direction.
The Hilbert--Schmidt Independence Criterion (HSIC) is a kernel measure of
statistical dependence that is zero iff two variables are independent and, unlike
correlation, detects \emph{non-linear} association without assuming a functional
form. Running HSIC on every pair of metrics reveals the linked chain and isolates
the unrelated ones.

\section{HSIC}

For paired samples $\{(a_i,b_i)\}_{i=1}^n$, with RBF kernels
$k(a,a')=\exp(-\|a-a'\|^2/\sigma_a^2)$ and $l(b,b')=\exp(-\|b-b'\|^2/\sigma_b^2)$
(bandwidths by the median heuristic), form Gram matrices $K,L\in\R^{n\times n}$
and the centering matrix $H=I_n-\tfrac1n\mathbf{1}\mathbf{1}^\top$. The empirical
HSIC is
\begin{equation}
\widehat{\HSIC}(a,b)=\frac{1}{(n-1)^2}\,\operatorname{tr}\!\big(KHLH\big)\;\ge 0,
\end{equation}
large when $a,b$ are dependent and $\approx 0$ when independent. It is symmetric
and makes \emph{no} directional / causal claim --- it is purely an association.

\section{Method}

\begin{enumerate}
\item Compute each metric per entity (a common index), giving columns
$m_1,\dots,m_p$.
\item For every pair compute $\widehat{\HSIC}(m_i,m_j)$.
\item Keep pairs with $\widehat{\HSIC}$ above a small threshold as
\emph{linkage edges}; the connected component is the metric \emph{chain}.
\end{enumerate}

\section{Result}

On a synthetic marketplace with a linked chain
\texttt{delivery\_mins} -- \texttt{user\_rating} -- \texttt{n\_orders} --
\texttt{gmv} (metrics across the order, user, and merchant levels) plus two
unrelated metrics (\texttt{quantity}, \texttt{discount}), the pairwise HSIC is:

\begin{table}[h]
\centering
\begin{tabular}{llr}
\toprule
metric $a$ & metric $b$ & $\widehat{\HSIC}$ \\
\midrule
delivery\_mins & user\_rating & 0.055 \\
delivery\_mins & n\_orders    & 0.056 \\
delivery\_mins & gmv          & 0.046 \\
user\_rating   & n\_orders    & 0.062 \\
user\_rating   & gmv          & 0.054 \\
n\_orders      & gmv          & \textbf{0.070} \\
\midrule
\multicolumn{2}{l}{\emph{any pair involving} \texttt{quantity} \emph{or} \texttt{discount}}
 & $\le 0.002$ \\
\bottomrule
\end{tabular}
\caption{Pairwise HSIC. The four chain metrics are mutually dependent
($\widehat{\HSIC}\in[0.046,0.070]$); every pair touching \texttt{quantity} or
\texttt{discount} is $\le 0.002$ (independent). HSIC recovers the linkage chain
\texttt{delivery\_mins}--\texttt{user\_rating}--\texttt{n\_orders}--\texttt{gmv}
and isolates the unrelated metrics --- with no explicit relationship supplied and
no causal direction claimed.}
\end{table}

\paragraph{Takeaway.} HSIC alone --- one symmetric, assumption-light, kernel
dependence number per pair --- is enough to \emph{discover the metric linkage
across levels}. Direction and root cause are deliberately \emph{not} claimed
here; this prototype only recovers which metrics are associated.

\end{document}
